\section{Discussione e problemi aperti}
\label{sec:discussion}

\subsection{Cosa è noto e cosa aggiungiamo}

La Tabella~\ref{tab:known} separa i risultati noti dalle osservazioni di questo lavoro.

\begin{table}[t]
\centering
\small
\begin{tabular}{@{}p{0.55\linewidth}p{0.38\linewidth}@{}}
\toprule
Enunciato & Fonte \\
\midrule
Potenziale 1RSB $\Xi(\beta,m)$, limiti entropico ed energetico & \citep{monasson1995structural,mezard2003cavity,mezard2002random} \\
SP, SP-$y$ e funzionale di complessità & \citep{mezard2002analytic,mezard2002random,braunstein2005survey} \\
SP come BP su uno spazio esteso, famiglia SP$(\rho)$ & \citep{braunstein2004surveybp,maneva2007survey} \\
Bias $b_i$ come fattore sul numero di cluster; BSP & \citep{marino2016bsp} \\
SP-$y$ con backtracking e indice con segno & \citep{battaglia2004minimizing} \\
Semianelli termodinamici, Shannon, Rényi, Tsallis, KL & \citep{marcolli2014thermodynamic} \\
Ensemble canonici generalizzati & \citep{costeniuc2005generalized,costeniuc2006generalized} \\
\midrule
Stazionarietà del funzionale completo, non della forma compatta (Lemma~\ref{lem:stationary}) & dimostrazione diretta nostra di un fatto generale~\citep{yedidia2005constructing} \\
Due livelli della somma di Shannon; $y$ locale uguale a $y$ dei cluster & nostro \\
Settori di SP-$y$ come somme con entropia relativa & nostro \\
Ordine di Rényi uguale a $m$; spettro multifrattale; $h_m=0$ per $m\ge m^*$ & nostro \\
Forma chiusa e limite $T\to0$ della somma di Tsallis deformata & nostro, con~\citep{curado1991generalized} \\
Tsallis sui cluster: tre regimi, $y_{\mathrm{eff}}$ & nostro \\
Tilt $\gamma$: rottura dell'invariante di MMW & nostro, da~\citep{maneva2007survey} \\
Identità locale con segno e resto di secondo ordine & nostro; il caso del valore più probabile è in~\citep{marino2016bsp} \\
Equazioni di Bellman soft e politica di Gibbs & \citep{haarnoja2018soft} \\
Look-ahead soft a due passi, valore a campi congelati, ricompensa di Tsallis (Sezione~\ref{sec:softq}) & nostro; verificato in Lean~4 \\
\bottomrule
\end{tabular}
\caption{Risultati noti (sopra) e osservazioni di questo lavoro (sotto).}
\label{tab:known}
\end{table}

Il quadro dei semianelli non produce da solo un algoritmo nuovo.
A $T$ fisso il semianello di Shannon è isomorfo al semianello di BP, quindi ogni equazione di cavità si può scrivere in esso.
Il valore del quadro sta nei limiti e nelle deformazioni.
Esso dice quali deformazioni restano associative, quali degenerano, e a quale livello della teoria agisce ogni parametro.

\subsection{Problemi aperti}

\paragraph{Non concavità di $\Sig_e$.}
La deformazione di Tsallis in scala $1/N$ con $q<1$ dà stati nuovi solo se $\Sig_e(e)$ ha tratti non concavi.
Per saperlo si può calcolare $\Sig_e(e)$ con la dinamica di popolazione vicino ad $\alpha_s$, per $K=3$ e $K=4$.

\paragraph{Il ciclo su $y_{\mathrm{eff}}$.}
Il Corollario~\ref{cor:yeff} suggerisce un ciclo esterno che aggiorna $y_{\mathrm{eff}}$ con l'energia del funzionale di Bethe.
Resta da capire se il ciclo converge e se, con backtracking, cambia la soglia algoritmica.

\paragraph{Il tilt $\gamma$.}
BP sul modello di MMW con $\omega_o=0$ e $\omega_*=e^\gamma$ è ben definita, ma non si chiude su una survey.
Resta da confrontare questa misura con la chiusura approssimata del nostro codice.
Resta anche da capire se la misura aiuta la ricerca di cluster non congelati, come suggerito in~\citep{marino2016bsp}.

\paragraph{Il resto $R_{k,s}$ su grafi con cicli.}
La Proposizione~\ref{prop:remainder} dà l'ordine del resto ma non la sua grandezza.
Una stima tramite la suscettività di SP e la lunghezza dei cicli renderebbe quantitativo il confronto fra $I_k$, $b_k$ e $b_{\mathrm{back}}$.
Ricci-Tersenghi e Semerjian~\citep{ricci2009cavity} studiano la decimazione con il metodo della cavità, e il loro approccio è un punto di partenza.

\paragraph{Esperimenti.}
Le prove numeriche fatte finora sono troppo piccole.
Un confronto utile deve fissare il budget di calcolo, usare molte istanze per ogni $(K,\alpha,N)$, e riportare la frazione di istanze risolte con il suo errore.
Le regole di rilascio da confrontare sono $b_k$, $b_{\mathrm{back}}$ e $I_k$, a parità di $r$ e $f$.

\section{Conclusioni}

Abbiamo scritto le deformazioni note di SP nella lingua dei semianelli termodinamici e abbiamo derivato ogni passo.
La somma di Shannon agisce su due livelli, e il livello dei cluster fissa il reweighting locale.
L'ordine di Rényi dei pesi dei cluster è il parametro di Parisi, e vicino alla soglia la parte utile dello spettro multifrattale si riduce a SP.
La deformazione di Tsallis con $q$ fisso degenera, mentre la deformazione in scala $1/N$ è un ensemble canonico generalizzato che si riduce a SP-$y$ con un $y$ efficace.
L'identità locale della complessità giustifica il bias di BSP per il fissaggio e dà una regola con segno per il rilascio, con un resto di secondo ordine.
I diagrammi dei limiti raccolgono questi fatti e mostrano quali varianti sono nuove e quali sono riparametrizzazioni di algoritmi noti.
